\section{Self Attention 的重要性质}

\subsection{置换等变性}

Self Attention 有一个重要的数学性质：\textbf{置换等变性}（Permutation Equivariance）。

\begin{figure}[H]
\centering
\includegraphics[width=0.95\textwidth]{slides-images/slide-034.jpg}
\caption{置换等变性：如果打乱输入向量的顺序，输出向量也会以相同方式打乱，但每个输出向量的值不变\protect\footnotemark}
\end{figure}
\footnotetext{Slides 第35页。}

\begin{importantbox}{Self Attention 操作于集合而非序列}
因为 Self Attention 具有置换等变性，它实际上\textbf{不关心输入向量的顺序}——它将输入视为一个无序\textbf{集合}（set），而非有序\textbf{序列}（sequence）。

这既是优势也是挑战：
\begin{itemize}
    \item \textbf{优势}：天然适合处理无序数据（如点云、图的节点集合）
    \item \textbf{挑战}：对于语言等有序数据，需要额外机制来注入位置信息
\end{itemize}
\end{importantbox}

\subsection{位置编码（Positional Encoding）}

为了让 Self Attention 感知输入的顺序，我们在每个输入向量上\textbf{拼接}（或相加）一个\textbf{位置编码}（Positional Embedding），告诉模型每个向量在序列中的位置。

\begin{figure}[H]
\centering
\includegraphics[width=0.95\textwidth]{slides-images/slide-035.jpg}
\caption{位置编码：在每个输入向量上附加位置信息，使 Self Attention 能区分不同位置的相同内容\protect\footnotemark}
\end{figure}
\footnotetext{Slides 第36页。}

\begin{knowledgebox}{位置编码的常见方法}
\begin{itemize}
    \item \textbf{正弦/余弦编码}（原始 Transformer）：使用不同频率的正弦函数，无需训练
    \item \textbf{可学习位置编码}：每个位置对应一个可训练的向量
    \item \textbf{旋转位置编码}（RoPE）：通过旋转矩阵将位置信息编码到 Query/Key 中，支持外推到更长序列
    \item \textbf{相对位置编码}：编码向量之间的相对距离而非绝对位置
\end{itemize}
\end{knowledgebox}

\subsection{本章小结}

Self Attention 天然具有置换等变性，本质上处理的是无序集合。位置编码是让注意力机制处理有序序列数据的关键补丁。

%% ========== Masked Attention ==========

